\subsection{Evaluator conventions for local moves}
\label{subsec:local-kirby-models}

The following conventions turn a topological Kirby move into a local ribbon
calculation while keeping the exterior tangle fixed.

\begin{convention}[Auxiliary orientation reversal]
\label{conv:orientation-reversal}
Fix a twisted categorical input \((\cC,\eta)\).  Let \(\mathcal H\) be a
one-ended Kirby presentation, let
\(a\in Z^2(\mathcal H)\), and suppose its \(k\)th framed component has
oriented generator \(e_k\), cocycle coordinate \(a_k\), and simple label
\(i\in\operatorname{Irr}(\cC)\).  Reversing the chosen orientation of that
framed \(2\)-handle
simultaneously sends
\[
  e_k\longmapsto-e_k,
  \qquad a_k\longmapsto-a_k,
  \qquad i\longmapsto i^*.
\]
Reversing a dotted \(1\)-handle orientation reverses its cross-section
coorientation and replaces its boundary word by the reversed rigid dual.  Both
operations change only auxiliary oriented bases, not the underlying Kirby
presentation.
\end{convention}

\begin{convention}[Coordinates for a \(2/2\) slide]
\label{conv:two-two-coordinates}
Fix a twisted categorical input \((\cC,\eta)\).
Let \(e_m\) and \(e_s\) be the moving and stationary \(2\)-handle generators,
let \(a_m,a_s\in\mathbb Z\) be the corresponding cocycle coordinates, and
let \(x,y\in\operatorname{Irr}(\cC)\) be their simple labels.  For the positive
slide we use
\[
  e_m'=e_m+e_s,
  \qquad (a_m',a_s')=(a_m+a_s,a_s),
\]
with fusion channels \(z\subset x\otimes y^*\).  For the negative slide we
use
\[
  e_m'=e_m-e_s,
  \qquad (a_m',a_s')=(a_m-a_s,a_s),
\]
with channels \(z\subset x\otimes y\).  These are the handle-basis change and
its dual cochain coordinates.
\end{convention}
